% Official CC BY4.0 published HTML A1; original math alttext and author text, literal original numbering.


\par\medskip\noindent\textbf{Appendix A Proofs on algebraic natural proofs}\par




In this section we prove Theorem 2.2.





\par\medskip\noindent\textbf{Definition A.1.}\par



We call a sequence $r\colon\mathbb{N}\to\mathbb{N}$ strongly superpolynomial if $r(n)$ dominates every polynomial sequence eventually, that is, for every polynomial $p$ there exists $N$ such that $r(n)>p(n)$ for all $n\geq N$.






\par\medskip\noindent\textbf{Lemma A.2.}\par



$(\Delta_{n})\in I(\mathcal{C})$ if and only if there exists a strongly superpolynomial sequence $r$ such that $\Delta_{n}(X_{n,r(n)})=0$.






\par\medskip\noindent\textbf{Proof.}\par




Clearly, if $\Delta_{n}(X_{n,r(n)})=0$, then $\Delta_{n}(X_{n,p(n)})$ vanishes eventually for every polynomially bounded sequence $p$, since $p(n)\leq r(n)$ eventually.





On the other hand, if $(\Delta_{n})$ is a sequence of metapolynomials in $I(\mathcal{C})$, we can define $r(n)$ as $r(n):=\max\{r\mid\Delta_{n}(X_{n,r})=0\}$.
For every polynomially bounded sequence $p(n)$ we have that $\Delta_{n}(X_{n,p(n)})$ vanishes eventually, so $r(n)\geq p(n)$ eventually.
$\hfill\blacktriangleleft$






\par\medskip\noindent\textbf{Proof of Theorem 2.2.}\par




We first observe we may assume $\deg(f_{n})=n$. Indeed, for every p-family $(f_{n})$ with the property that $\deg(f_{n})$ is strictly increasing, let $(f^{\prime}_{n})$ be the sequence defined by



\[\begin{aligned}
f^{\prime}_{m}=\left\{\begin{array}[]{ll}f_{n}&\text{if $m=\deg(f_{n})$ for some $n$}\\ 0&\text{otherwise}\end{array}\right.
\end{aligned}\]


Since $(f_{n})$ is a p-family, the sequence $\deg(f_{n})$ is polynomially bounded, hence $(f^{\prime}_{n})$ is a p-family as well. Moreover $(f_{n})\in\mathcal{C}$ if and only if $(f_{n})\in\mathcal{C}$, and similarly for $\bar{\mathcal{C}}$. Finally, for a sequence of metapolynomials $(\Delta_{n})$ of format $(*,n,*)$, we have $\Delta_{\deg(f_{n})}(f_{n})=0$ eventually if and only if $\Delta_{n}(f^{\prime}_{n})=0$ eventually. In particular, proving the statement for $(f^{\prime}_{n})$ is equivalent to proving the statement for $(f_{n})$. Since $\deg(f^{\prime}_{n})=n$, this shows we may assume $\deg(f_{n})=n$.





Suppose $(f_{n})\in\overline{\mathcal{C}}$, that is, the sequence $p(n)=\underline{c}(f_{n})$ is polynomially bounded.
For every sequence of metapolynomials $(\Delta_{n})\in I(\mathcal{C})$ we have that $\Delta_{n}(X_{n,p(n)})$ vanishes eventually.
Since the metapolynomials $\Delta_{n}$ are continuous, $\Delta_{n}(\overline{X_{n,p(n)}})$ vanishes eventually, so $\Delta_{n}(f_{n})=0$ eventually.





Conversely, suppose $(f_{n})\notin\overline{\mathcal{C}}$. Therefore $r(n)=\underline{c}(f_{n})$ is not polynomially bounded. Let $\theta(n)=\log_{n}r(n)$, $\hat{\theta}(n)=\max\{\theta(m)\mid m\leq n\}$, and $R(n)=\lceil n^{\hat{\theta}(n)}\rceil$. Note that $R(n)$ is monotonically increasing.





The sequence $R(n)$ is strongly superpolynomial.
Indeed, since $r(n)$ is not polynomially bounded, for every $k$ there exists $n_{0}$ such that $\theta(n_{0})\geq k$ and, therefore, $\hat{\theta}(n)\geq k$ for all $n\geq n_{0}$.
It follows that for every $k$ we have $R(n)\geq n^{k}$ eventually.





For all $\Theta$, if $n_{\Theta}$ is the minimal value such that $\theta(n_{\Theta})\geq\Theta$, then $\hat{\theta}(n_{\Theta})=\theta(n_{\Theta})$.
Since $\theta$ is unbounded, it follows that $\theta$ and $\hat{\theta}$ coincide on an infinite subset of $\mathbb{N}$.
On this subset we also have $R(n)=r(n)$.





Consider a metapolynomial sequence $\Delta_{n}$ such that $\Delta_{n}(\overline{X_{n,R(n)-1}})=0$, and if $r(n)=R(n)$, then $\Delta_{n}$ is chosen so that $\Delta_{n}(f_{n})\neq 0$ (this is possible because $\underline{c}(f_{n})=r(n)=R(n)>R(n)-1$).
Since $R(n)$ is strongly superpolynomial, $R(n)-1$ is strongly superpolynomial, and $(\Delta_{n})\in I(\mathcal{C})$ by Lemma A.2, and since $R(n)=r(n)$ infinitely often, $\Delta_{n}(f_{n})\neq 0$ infinitely often.
$\hfill\blacktriangleleft$




